% =============================================
% 实验报告 / 课程报告 LaTeX 通用模板
% 基于 xelatex + ctex 编译
% 适用：A4 / 12pt / 中文（宋体 + 黑体）
% =============================================
\documentclass[12pt, a4paper]{article}

% ---- 中文支持 ----
% 注：以下字体路径基于 macOS。若在其他系统上使用，请把 Path=... 替换为本机字体目录
% 或直接删除 Path 与 FontIndex 选项，改用系统已注册的字体名称（如 SimSun、SimHei）。
\usepackage[fontset=none]{ctex}
\setCJKmainfont{Songti.ttc}[Path=/System/Library/Fonts/Supplemental/, FontIndex=3, BoldFont=Songti.ttc, BoldFeatures={FontIndex=3}]
\setCJKsansfont{STHeiti Light.ttc}[Path=/System/Library/Fonts/, BoldFont=STHeiti Light.ttc]
\setCJKmonofont{Songti.ttc}[Path=/System/Library/Fonts/Supplemental/, FontIndex=3, BoldFont=Songti.ttc, BoldFeatures={FontIndex=3}]
\setCJKfamilyfont{zhsong}{Songti.ttc}[Path=/System/Library/Fonts/Supplemental/, FontIndex=3, BoldFont=Songti.ttc, BoldFeatures={FontIndex=3}]
\setCJKfamilyfont{zhhei}{STHeiti Light.ttc}[Path=/System/Library/Fonts/, BoldFont=STHeiti Light.ttc]
\setCJKfamilyfont{zhfs}{Songti.ttc}[Path=/System/Library/Fonts/Supplemental/, FontIndex=3, BoldFont=Songti.ttc, BoldFeatures={FontIndex=3}]
\setCJKfamilyfont{zhkai}{Songti.ttc}[Path=/System/Library/Fonts/Supplemental/, FontIndex=3, BoldFont=Songti.ttc, BoldFeatures={FontIndex=3}]
\newcommand{\songti}{\CJKfamily{zhsong}}
\newcommand{\heiti}{\CJKfamily{zhhei}}
\newcommand{\fangsong}{\CJKfamily{zhfs}}
\newcommand{\kaishu}{\CJKfamily{zhkai}}

% ---- 页面布局 ----
\usepackage[
  top=2.5cm, bottom=2.5cm,
  left=3cm,  right=2.5cm
]{geometry}

% ---- 图片 ----
\usepackage{graphicx}
\usepackage{float}
\graphicspath{{figures/}}

% ---- 表格 ----
\usepackage{booktabs}
\usepackage{tabularx}
\usepackage{multirow}

% ---- 代码高亮 ----
\usepackage{listings}
\usepackage{xcolor}

\lstset{
  basicstyle=\ttfamily\small,
  keywordstyle=\color{blue}\bfseries,
  commentstyle=\color{gray}\itshape,
  stringstyle=\color{red!70!black},
  numbers=left,
  numberstyle=\tiny\color{gray},
  numbersep=8pt,
  breaklines=true,
  frame=single,
  rulecolor=\color{black!30},
  backgroundcolor=\color{gray!5},
  showstringspaces=false,
  tabsize=4,
}

% ---- 超链接 ----
\usepackage[
  colorlinks=true,
  linkcolor=blue!70!black,
  urlcolor=blue,
  citecolor=green!50!black
]{hyperref}

% ---- 数学 ----
\usepackage{amsmath, amssymb}

% ---- 页眉页脚 ----
\usepackage{fancyhdr}
\pagestyle{fancy}
\fancyhf{}
\fancyhead[L]{\small \courseName}
\fancyhead[R]{\small \leftmark}
\fancyfoot[C]{\thepage}
\renewcommand{\headrulewidth}{0.4pt}

% ---- 其他工具 ----
\usepackage{enumitem}
\usepackage{caption}
\captionsetup{font=small, labelfont=bf}

% ---- 封面绘图 ----
\usepackage{tikz}
\usepackage{array}

% =============================================
%  封面信息（每次新报告只需修改这一段）
% =============================================
\newcommand{\courseName}{课程名称}                   % 页眉左侧 / 课程名
\newcommand{\expTitle}{实验大标题}                   % 封面正中大字
\newcommand{\expFullTitle}{实验报告完整题目}         % 题目栏
\newcommand{\academicYear}{2025-2026-2}              % 学年学期
\newcommand{\college}{计算机学院}
\newcommand{\major}{计算机科学与技术}
\newcommand{\studentName}{姓名}
\newcommand{\studentID}{学号}
\newcommand{\instructor}{任课教师}
\newcommand{\expDate}{YYYY 年 M 月 D 日}

% 校名图片文件名（位于 figures/ 下，不带路径前缀）
\newcommand{\titleLogo}{杭电文字 title.png}

% =============================================
\begin{document}

% ---- 封面 ----
\begin{titlepage}
  \thispagestyle{empty}

  \centering
  \vspace*{2.0cm}

  % ---- 校名图片 ----
  \includegraphics[width=11cm]{\titleLogo}

  \vspace{2.2cm}

  % ---- 实验大标题 ----
  {\fontsize{38}{50}\selectfont\heiti\bfseries \expTitle\par}

  \vspace{0.9cm}

  % ---- 学年 ----
  {\Large\bfseries （\academicYear）\par}

  \vspace{3.0cm}

  % ---- 信息表格 ----
  % 辅助命令：带下划线的定宽内容框
  \newcommand{\infobox}[1]{\underline{\makebox[7.8cm][c]{\large #1}}}

  \renewcommand{\arraystretch}{1.75}
  \begin{tabular}{>{\large\bfseries}p{2.6cm} l}
    题\quad 目 & \infobox{\expFullTitle}   \\
    学\quad 院 & \infobox{\college}        \\
    专\quad 业 & \infobox{\major}          \\
    学\quad 号 & \infobox{\studentID}      \\
    学生姓名   & \infobox{\studentName}    \\
    任课教师   & \infobox{\instructor}     \\
    完成日期   & \infobox{\expDate}        \\
  \end{tabular}

\end{titlepage}

% ---- 目录 ----
\tableofcontents
\newpage

% =============================================
%  正文
%  小提示：
%   - 章节多时可用 \input{sections/01-purpose.tex} 等方式拆分
%   - 图片放 figures/ 下，\includegraphics 时省略路径前缀
%   - 代码块用 lstlisting 环境，已在 \lstset 中预设样式
% =============================================

\section{实验目的}

% TODO: 用 itemize 列出本次实验需要达成的目标
\begin{itemize}
  \item 目的一……
  \item 目的二……
  \item 目的三……
\end{itemize}


\section{实验背景与原理}

% TODO: 阐述与本次实验相关的背景知识、概念定义和原理。
% 建议分自然段叙述，必要时配合 itemize/enumerate 罗列要点。


\section{实验环境}

% TODO: 列出操作系统、虚拟机、发行版、内核版本等。
\begin{table}[H]
  \centering
  \caption{实验环境信息}
  \begin{tabularx}{\linewidth}{>{\centering\arraybackslash}p{3.5cm} >{\centering\arraybackslash}X}
    \toprule
    \textbf{项目} & \textbf{说明} \\
    \midrule
    操作系统      & TODO \\
    虚拟机软件    & TODO \\
    Linux 发行版  & TODO \\
    内核版本      & TODO \\
    \bottomrule
  \end{tabularx}
\end{table}


\section{实验设计}

% TODO: 说明实验整体思路与各部分的安排。可以分若干 \subsection 介绍每个部分。

\subsection{第一部分：标题}

% TODO: 描述本部分的目的与所用命令/方法。

\subsection{第二部分：标题}

% TODO: 同上。


\section{实验过程与结果分析}

\subsection{第一部分：标题}

% TODO: 在每个步骤里穿插 \begin{lstlisting} 代码块和 \begin{figure} 截图。

\begin{enumerate}[label=\arabic*.]
  \item \textbf{操作步骤一。}

\begin{lstlisting}[language=bash, caption={示例命令}]
echo "hello world"
\end{lstlisting}

        在此处分析图 \ref{fig:example} 的输出：……

        \begin{figure}[H]
          \centering
          \includegraphics[width=0.9\linewidth]{example.png}
          \caption{示例截图}
          \label{fig:example}
        \end{figure}

  \item \textbf{操作步骤二。} ……
\end{enumerate}


\section{实验总结}

% TODO: 总结主要结论、收获与可拓展的方向。


\newpage
\section{参考文献}

\begin{enumerate}[label={[\arabic*]}, leftmargin=2.2em]
  \item 作者. \emph{文献标题}. 出版方, 年份. \url{https://example.com/}.
\end{enumerate}

\end{document}